\section*{September 7, 2026: Joint choices of size and organization}
{\small Drafted by Codex from the supplied framework and inspected repository.}

Jacob requested that the supplied September framework replace the earlier
notch-and-kink note. The new framework distinguishes a reduction in an economic
parent's proposed units from a change in how those units are divided across
constituents. It allows ordinary historical multi-constituent development and
mixed configurations in which only some constituents face the modeled burden.
The relevant legal mapping remains a maintained assumption to check against
administrative evidence.

The proposed counterfactual retains historical size and organization jointly,
using site characteristics to reweight the recent historical population. Its
primary quantity is the difference in proposed parent units conditional on
filing. The framework does not identify entry effects, completed housing, or
welfare. The excess at 99 alone no longer identifies a unique source frontier
once organization is a choice.

The current code constructs parent cohorts, constituent vectors, exposure
classifications, predetermined site characteristics, and descriptive and
reweighted distributions. It does not yet estimate the joint model. The
supplied framework's implementation sequence is a proposal for subsequent work;
this replacement changes no sample, linkage decision, or estimator. Comparable
follow-up windows, feasible legal partitions, and the treatment of parents with
more than two constituents remain issues for that implementation.

The replacement source is preserved verbatim. Its opening empirical diagnostics
describe an earlier figure-guide snapshot; they differ from the calibration,
local-moment, and splitting-verification outputs inspected in the current
working tree. Those diagnostics were not refreshed during this document
replacement. The current empirical outputs should be used when reporting the
latest descriptive results.

{\small\textit{Reproduction note.}
The supplied source was
\texttt{nyc\_size\_and\_splitting\_framework.tex}, accompanied by its PDF.
The canonical source is now \texttt{framework\_writeup.tex} at the repository
root; \texttt{make framework-writeup} compiles its PDF using its inline
bibliography. This entry records the working tree on branch
\texttt{simplify\_empirical\_tasks}, based on commit \texttt{7bbf6d5}, including
pre-existing uncommitted empirical changes. The source horizon for recent
filings is July 8, 2026. The document build is checked separately from the
empirical pipeline, which was read but not rerun for this replacement.}

\clearpage
\section*{September 7, 2026: Revision around the housing-supply question}
{\small Drafted by Codex after Jacob requested a clearer framework.}

The revised note leads with the distinction between fewer proposed apartments
and a different division of the same apartments across constituents. It moves
from that question to the historical comparison, the developer's joint choice,
and the interpretation of costs. Data construction, partition formulas, and
numerical checks now sit in the appendices.

Jacob asked to remove the optional historical model. The revision retains
observed historical size and organization directly and removes the historical
propensity regression and its logistic probability formulas. It preserves
heterogeneous ordinary organization costs. Their conditional magnitude
distribution remains a structural specification to choose; the historical
choice restricts preferences but cannot identify their strength. The note
states this limitation without adding a replacement parametric model.

The empirical table and sample counts now read the same generated values as
the current figure guide. This supersedes the earlier imported diagnostics.
The update changes the reported sample size, local differences, and
splitting-verification counts to match the saved outputs; it changes no
underlying analytical result. The framework distinguishes the complete-site
sample used for reweighting from the broader sample of observed 198-unit
parents, and explains the potential selection at the 50-unit sample boundary.

Table 1 now displays the weighted historical and observed post-period shares
alongside each difference and confidence interval. The table note states the
subtraction direction and explains that differences use unrounded shares.
Independent recalculation from parent records reproduced the point estimates
and calibration weights; the saved bootstrap draws reproduced the intervals.

{\small\textit{Reproduction note.}
Run \texttt{make framework-writeup} at the repository root. The build checks
\path{build_figure_guide_values.R} and its upstream task dependencies,
then compiles the framework. Run \texttt{make} in \texttt{logbook/} for this
entry. Source definitions, empirical estimators, and linkage decisions are unchanged.
The value-export script now also supplies the shares shown in Table 1.}
